% =========================================================================
% Chapter 8 - The role of Model 2
% =========================================================================
\chapter{The role of Model 2: a learned kinetic energy functional}
\label{ch:model2}

\section{The question, stated without charity}

A skeptic's version of the objection is worth writing down, because the weak
answers to it are the ones most often given:

\begin{quote}
\itshape
You already have Model~1, which predicts $\rho$ from the geometry. If you have
a converged DFT calculation you have $\tau$ for free; and if you do not have
one, you have no $\rho$ to feed Model~2 either. What is \opt{chg2tau}
\emph{for}?
\end{quote}

Three answers are available. Two are weak.

\begin{description}
  \item[``$\tau$ is a useful observable'' --- weak.] It enters meta-GGA
    functionals, the electron localisation function and bonding analysis. True,
    but anyone holding a converged \file{CHGCAR} also holds a \file{TAUCAR}.
    Predicting a quantity one already has is a benchmark, not an application.
  \item[``It completes the pipeline'' --- weak.] Chaining
    $\vext\to\rho\to\tau$ yields $\tau$ for a structure never computed, which
    is useful for screening. But it makes Model~2 a convenience: nothing in it
    is more than interpolation, and its errors feed back into nothing.
  \item[``It \emph{is} the orbital-free kinetic energy functional'' --- the
    actual reason.] See below.
\end{description}

$\Ts[\rho]$ is the \emph{only} term of the total energy with no accurate
explicit density functional. Kohn--Sham theory evades the problem by
reintroducing orbitals and pays $\mathcal{O}(N^3)$ for it. Orbital-free DFT
removes them and scales near-linearly, but is only as good as its functional.
A model of $\rho\mapsto\tau$ \textbf{is} such a functional, because
\begin{equation}
  \Ts[\rho] = \int \tau[\rho](\rr)\,\mathrm{d}\rr .
\end{equation}
Model~2 is therefore not a predictor of a convenient observable; it is a
candidate solution to the central open problem of orbital-free DFT, in the one
form a network can represent and autograd can differentiate. \textbf{Model~2 is
the scientific contribution; Model~1 is the infrastructure around it.}

\section{Is \texorpdfstring{$\tau$}{tau} a functional of \texorpdfstring{$\rho$}{rho}?}

Worth settling, because the answer is subtler than ``Hohenberg--Kohn says so''.

For the integrated quantity the argument is clean: HK gives $\rho\to
v_{\mathrm{s}}$ uniquely for non-interacting $v$-representable densities,
$v_{\mathrm{s}}$ determines the orbitals, and the orbitals determine $\Ts$. The
same chain $\rho\to v_{\mathrm{s}}\to\{\psi_i\}\to\tau$ establishes
$\tau[\rho]$ as well. But note what that chain contains --- solving the
Kohn--Sham equations. The functional is thus

\begin{itemize}
  \item \textbf{exact but non-local}: $\tau$ at one point depends on $\rho$
    everywhere, because the orbitals are delocalised Bloch states;
  \item \textbf{not semi-local}: no finite gradient expansion converges to it,
    which is precisely why Thomas--Fermi and its gradient corrections fail;
  \item \textbf{expensive to evaluate exactly}: the definition requires the very
    diagonalisation orbital-free DFT exists to avoid.
\end{itemize}

Two architectural consequences follow. The architecture \emph{must} be
non-local, which rules out a finite-receptive-field convolutional network and
motivates the globally coupled Fourier layer of Chapter~\ref{ch:fno}. And the
target is well posed: unlike most machine-learning targets, $\tau[\rho]$ is a
genuine single-valued functional, so failure to learn it is a capacity or data
problem, not an ill-posedness problem.

\begin{pwarn}
$\tau$ is a functional of $\rho$ \emph{for ground-state, non-interacting
$v$-representable densities}. Inside a self-consistency loop the intermediate
densities are not ground states of any $v_{\mathrm{s}}$, so strictly
$\tau[\rho]$ is undefined there. Every orbital-free calculation uses the
functional anyway --- but it means a model trained only on converged densities
is \emph{extrapolated} every time it is used variationally. This is the single
largest risk in the programme.
\end{pwarn}

\section{What the model must get right --- and it is not \texorpdfstring{$\tau$}{tau}}
\label{sec:notthetau}

Orbital-free DFT does not consume $\tau$. It consumes the functional
derivative
\begin{equation}
  v^{\mathrm{kin}}_{\mathrm{s}}(\rr) \equiv \frac{\delta \Ts}{\delta\rho(\rr)},
\end{equation}
which is the term appearing in Eq.~\eqref{eq:euler} and therefore the term that
determines the density.

A model can have small $\lVert\tau_\theta-\tau\rVert$ and a \emph{useless}
derivative, because differentiation amplifies high-frequency error: a ripple of
amplitude $\epsilon$ and wavevector $G$ contributes $\epsilon$ to the value but
$\epsilon G$ to the gradient. Three consequences:

\begin{enumerate}
  \item Reporting only pointwise $\tau$ error reports a \textbf{proxy}. The
    relative $L^2$ figures of Chapter~\ref{ch:results} are necessary, not
    sufficient.
  \item The $H^1$ data loss exists for this reason: it penalises the gradient
    error that plain $L^2$ ignores.
  \item The definitive test is \textbf{downstream} --- insert the functional
    into an orbital-free minimiser and measure the converged density and energy.
\end{enumerate}

\section{The mechanism: \texorpdfstring{$\delta \Ts/\delta\rho$}{dTs/drho} by autograd}
\label{sec:autograd}

This is what makes a \emph{neural} functional more useful than a tabulated one,
and it is mechanically simple: $\Ts$ is a scalar computed from $\rho$ by a
differentiable program, so one backward pass yields
$\partial \Ts/\partial\rho_i$ at every grid point --- no finite differences and
no derivative to derive by hand.

\begin{lstlisting}
rho = rho_physical.detach().requires_grad_(True)      # (B,1,Nx,Ny,Nz)

tau = target_transform.inverse(model2(input_transform(rho), cell))
T_s = integrate(tau, cell)                            # (B,), eV

grad, = torch.autograd.grad(T_s.sum(), rho, create_graph=True)
dTs_drho = grad / volume_element                      # see below
\end{lstlisting}

\subsection{The discretisation factor}

Take $\Ts = \sum_i \tau_i \Delta v$. The functional derivative is defined by
\[
  \delta \Ts = \int \frac{\delta \Ts}{\delta\rho}\,\delta\rho\,\mathrm{d}\rr
  \approx \sum_i \left(\frac{\delta \Ts}{\delta\rho}\right)_i
          \delta\rho_i\,\Delta v ,
\]
while autograd returns
$\delta \Ts = \sum_i (\partial \Ts/\partial\rho_i)\,\delta\rho_i$. Hence
\begin{equation}
  \boxed{\;
    \frac{\delta \Ts}{\delta\rho}(\rr_i)
      = \frac{1}{\Delta v}\,\frac{\partial \Ts}{\partial \rho_i} \;}
  \qquad \Delta v = \frac{\Omega}{N_1N_2N_3} .
\end{equation}

\begin{pwarn}
Omitting the $1/\Delta v$ rescales the entire kinetic potential by the number of
grid points --- a factor of $3\times10^{4}$ on these grids --- which would
silently destroy any Euler--Lagrange residual built from it, without raising
anything. \opt{create\_graph=True} is likewise required if the derivative is to
appear in a loss that is itself backpropagated.
\end{pwarn}

\subsection{Validation of the derivative}

The derivative is meaningful only if verified independently:

\begin{center}
\begin{tabular}{@{}ll@{}}
\toprule
Test & Expectation \\
\midrule
Substitute analytic Thomas--Fermi for Model 2
  & recovers $\tfrac{5}{3}C_{\mathrm{TF}}\rho^{2/3}$ \\
Substitute von Weizsäcker
  & recovers $-\tfrac12\nabla^2\sqrt\rho/\sqrt\rho$ \\
Finite-difference a few random voxels & agrees to FD accuracy \\
Uniform density & $\delta \Ts/\delta\rho$ constant in space \\
\bottomrule
\end{tabular}
\end{center}

These are cheap and they catch the $\Delta v$ error, a broken transform chain,
and any accidental \opt{detach()}.

\section{Coupling: how Model 2 trains Model 1}

\subsection{The equation}

The Euler--Lagrange condition~\eqref{eq:euler} holds pointwise at the ground
state with $\mu$ constant in space. Every term is available:

\begin{center}
\begin{tabular}{@{}lll@{}}
\toprule
Term & Source & Exact? \\
\midrule
$\vext$ & Model~1's \emph{input} & exact (tabulated pseudopotential) \\
$\rho$ & Model~1's \emph{output} & learned \\
$\vH[\rho]$ & $4\pi e^2\rho(\GG)/G^2$, one FFT & \textbf{exact} \\
$\delta \Ts/\delta\rho$ & autograd through Model~2 & learned \\
$\vxc[\rho]$ & LDA/GGA, or omitted & approximate \\
$\mu$ & eliminated by the cell average & --- \\
\bottomrule
\end{tabular}
\end{center}

Subtracting the cell average removes the unknown $\mu$ exactly, leaving the
residual of Eq.~\eqref{eq:elresidual}, which must vanish for the true
ground-state density.

\subsection{Why this is worth the trouble}

The residual is computable from \textbf{Model~1's own input and output}, plus
Model~2. It requires \textbf{no additional labels} --- no extra DFT
calculations, no targets --- so it can be evaluated on unlabelled structures,
on perturbed or interpolated geometries generated on the fly, and on the
model's own predictions during training.

That converts Model~2 from a passive predictor into an \textbf{active physical
constraint on Model~1}. This coupling is what the architecture is built around.

\subsection{The loss}

\begin{equation}
  \mathcal{L}_1 =
    \underbrace{\bigl\lVert \rho_\theta - \rho^{\mathrm{DFT}} \bigr\rVert_{\mathrm{rel}}}_{\text{data}}
    + \lambda_{\mathrm{EL}}
      \underbrace{\Bigl\langle \bigl(r(\rr)/s\bigr)^2 \Bigr\rangle}_{\text{Euler--Lagrange}}
    + \lambda_{N}
      \underbrace{\Bigl(\tfrac{\int\rho_\theta - N}{N}\Bigr)^2}_{\text{particle number}}
\end{equation}

with $s=\operatorname{std}(\vext)$ rendering the residual dimensionless and
comparable across materials whose potentials differ by an order of magnitude.
Three rules apply, each established elsewhere in this project: physics terms act
on the prediction decoded to \emph{physical} units; every term is
dimensionless; and $\lambda_{\mathrm{EL}}$ is ramped only after the data term
has converged, since a residual evaluated at random initialisation is noise.

\subsection{Freeze Model 2, or train jointly?}

\begin{description}
  \item[Frozen.] Model~2 is a fixed physics oracle. Simple, stable, and the
    constraint is exactly as good as the functional. Recommended first.
  \item[Joint.] Mutual regularisation, and Model~2 is exposed to the densities
    Model~1 actually produces --- which directly attacks the distribution-shift
    problem. But there is a real hazard.
\end{description}

\begin{pwarn}
\textbf{The Euler--Lagrange residual alone has trivial solutions.} With both
models free, the pair can drive $r\to0$ by co-adapting: Model~2 can learn a
functional whose derivative cancels $\vext+\vH+\vxc$ for whatever $\rho$
Model~1 emits, without either being correct. The residual is a
\emph{consistency} condition, not a \emph{correctness} condition. Joint training
is safe only with both data terms retained and dominant; the diagnostic of
collapse is a residual that falls while the held-out data error rises.
\end{pwarn}

\section{Critical assessment}

\subsection{Genuine strengths}

\begin{itemize}
  \item The target is a well-posed functional, not a fitted correlation.
  \item It attacks the actual bottleneck of orbital-free DFT.
  \item The derivative is free and exact via autograd --- removing the
    traditional obstacle to proposing new kinetic functionals, which is
    deriving $\delta\Ts/\delta\rho$ by hand.
  \item Exact constraints exist and can be imposed structurally:
    $\tau\ge\tauvW$ already is, supplying $\sim31\,\%$ of $\tau$ analytically.
  \item On this data it beats every classical functional by roughly an order of
    magnitude, with Thomas--Fermi at negative $R^2$.
\end{itemize}

\subsection{Real weaknesses}

\begin{itemize}
  \item \textbf{Distribution shift} is the central risk, and it is a domain
    question rather than a practical inconvenience: the functional is strictly
    undefined off the $v$-representable manifold.
  \item \textbf{We optimise $\tau$ and use $\delta\Ts/\delta\rho$}
    (Section~\ref{sec:notthetau}). Until training runs through the derivative,
    the objective is a proxy.
  \item \textbf{The derivative is unvalidated.} Nothing in the results of
    Chapter~\ref{ch:results} bears on $\delta\Ts/\delta\rho$; the
    Section~\ref{sec:autograd} checks are not yet implemented.
  \item \textbf{Data scale.} Three structures of one element.
  \item \textbf{The Pauli potential constraint is unused.} Levy--Ou-Yang gives
    $v_{\mathrm{P}}=\delta T_{\mathrm{P}}/\delta\rho\ge0$ --- a constraint on
    the \emph{derivative}, which is exactly the object that should be
    constrained.
\end{itemize}

\section{Verification plan}

Falsifiable, in dependency order:

\begin{enumerate}
  \item \textbf{Derivative correctness} --- the substitution tests above.
    Blocking: nothing downstream is meaningful until they pass.
  \item \textbf{Residual floor} --- evaluate $r(\rr)$ with the \emph{reference}
    $\rho$ and a classical functional. It will not vanish, since the functional
    is approximate; record the floor, because a learned functional must beat it
    to be adding anything.
  \item \textbf{Ablation} --- $\lambda_{\mathrm{EL}}=0$ versus $>0$ on held-out
    material, matched otherwise. Same protocol as the Pauli-head comparison of
    Section~\ref{sec:headresult}.
  \item \textbf{Label efficiency} --- the real promise. Train Model~1 on $k$
    labelled structures plus $m$ unlabelled ones under the residual, and show
    the curve beats $k$ labelled alone. If the residual cannot buy label
    efficiency, its main practical argument fails.
  \item \textbf{Downstream} --- insert the functional into an orbital-free
    minimiser and measure converged density and energy error. The only test
    that matters.
\end{enumerate}

\section{Status}

\begin{center}
\begin{tabular}{@{}ll@{}}
\toprule
Component & State \\
\midrule
$\rho\mapsto\tau$ supervised & done (rel.\ $L^2$ $0.0620$, $R^2$ $0.9924$) \\
$\tau\ge\tauvW$ structural & done (\pyclass{PauliResidualOperator}) \\
Exact spectral $\nabla$, $\nabla^2$, $\vH$ & done \\
Euler--Lagrange residual, analytic functional & done \\
$\delta\Ts/\delta\rho$ via autograd & \textbf{done}
(\pyclass{kinetic\_potential}, \pyclass{operator\_kinetic\_potential}) \\
Euler--Lagrange residual, \emph{learned} functional & done in the library
(\opt{kinetic=}); not wired into training \\
Joint training & not implemented \\
Pauli potential $v_{\mathrm{P}}\ge0$ & not implemented \\
Orbital-free minimiser integration & not implemented \\
\bottomrule
\end{tabular}
\end{center}

The autograd derivative is what turns two regressions into a differentiable
orbital-free solver, and it is exact: against analytic Thomas--Fermi it
reproduces the closed-form potential to $1\times10^{-14}$~eV, and it satisfies
Euler's homogeneity sum rule
$\int (\delta\Ts/\delta\rho)\,\rho\,d^3r = \tfrac53 \Ts$ to ten decimal places.
On a rough density it is \emph{more} reliable than the hand-derived analytic
von Weizs\"acker potential, which drifts from that sum rule by several per cent
where the autograd route stays exact.

\begin{pwarn}[title={Accuracy in $\tau$ does not imply accuracy in its derivative}]
Measured on a model trained to reproduce analytic Thomas--Fermi: relative
$L^2$ of $3.0\times10^{-3}$ on $\tau$, and $9.2\times10^{-1}$ on
$\delta\Ts/\delta\rho$ --- a $307\times$ amplification, with the \emph{mean}
wrong by a factor of fifteen.

A mean that wrong is not high-frequency amplification, which would give noise
about the right value. It is a direction the training data never explored: the
derivative's mean is exactly the response of $\Ts$ to a \emph{uniform}
rescaling of $\rho$, and every training density had the same mean. Retrained on
densities whose mean varies, the same architecture with the same $\tau$ error
gave $1.6\times10^{-1}$, and the sum rule went from $0.057$ to $0.973$.

The sum rule is the diagnostic, and it is computable from a learned functional
with no reference potential --- the only check available when there is no
closed form. Read it before closing the loop.
\end{pwarn}


\section{Distilling the Pauli factor, and which candidate to quote}
\label{sec:pareto-knee}

What Model 2 has to get right is the Pauli enhancement factor
$F = (\tau - \tauvW)/\tauTF$, so that is what the symbolic search fits.
A search does not return \emph{an} expression: it returns a front, every
candidate that was the best of its length. The front states a trade and
refuses to resolve it.

The \textbf{knee} is where resolving it costs least --- the candidate nearest
the ideal corner $(0,0)$ once both axes are rescaled to $[0,1]$ across the
front:
\[
  d_i = \sqrt{\hat c_i^{\,2} + \hat{\mathcal{L}}_i^{\,2}},
  \qquad
  \hat x_i = \frac{x_i - \min_j x_j}{\max_j x_j - \min_j x_j} .
\]

Two choices in that definition are load-bearing.

\textbf{Both axes are rescaled.} A node count and a loss have no common unit,
so a Euclidean distance between them is meaningless until they are made
comparable. Rescaling across the front is what gives the metric a meaning
relative to the range the search actually produced.

\textbf{The loss enters logarithmically.} A front spans orders of magnitude. On
a linear scale every candidate but the most accurate collapses onto
$\hat{\mathcal{L}}\approx1$, and the knee degenerates to the shortest
expression whatever it cost --- which is not a trade-off, it is a constant.

The report carries both candidates: the lowest-loss expression and the knee,
each with its own parity plot on identical voxels, laid side by side. Where the
two clouds are indistinguishable the extra nodes bought nothing, and that is
visible rather than summarised.

\begin{pnote}
A knee is a heuristic, not a theorem. With one candidate it is that candidate;
where the front is a straight line every point is equidistant. It answers
\emph{which expression is worth its length}. Whether an expression is
\emph{right} is what the asymptotic limits test ---
and those are statements about physics the search was never shown.
\end{pnote}
